\begin{scanpage}{0055}{38}
\[\left.\begin{aligned}E_z=-H_\phi&=E_{0z}[f^0(1)e^{ika(\sin\alpha-\sin\phi)}+\\&\qquad+f^0(2)e^{-ika(\sin\alpha-\sin\phi)}]\frac{e^{i(kr+\frac\pi4)}}{\sqrt{2\pi kr}},\\E_\phi=H_z&=H_{0z}[g^0(1)e^{ika(\sin\alpha-\sin\phi)}+\\&\qquad+g^0(2)e^{-ika(\sin\alpha-\sin\phi)}]\frac{e^{i(kr+\frac\pi4)}}{\sqrt{2\pi kr}}.\end{aligned}\right\}\tag{6.08}\]
Here the first terms correspond to the waves from edge 1 $(y=a)$, and the second terms correspond to the waves from edge 2 $(y=-a)$. The functions $f^0$ and $g^0$ are determined in the right half-plane $\left(|\phi|<\frac\pi2\right)$ by the equations
\[\left.\begin{aligned}f^0(1)&=-f^0(2)=\frac{\cos\alpha}{\sin\alpha-\sin\phi},\\g^0(1)&=-g^0(2)=\frac{\cos\phi}{\sin\alpha-\sin\phi}.\end{aligned}\right\}\tag{6.09}\]
Now let us find the field radiated by the nonuniform part of the current. Assuming the strip is sufficiently wide $(ka\gg1)$, one is able to approximately consider that the current near its upper edge is the same as on the ideally conducting half-plane $-\infty<y<a$, and near the lower edge it is the same as on the half-plane $-a<y<\infty$. Therefore, in accordance with \S\ 4, the field from the nonuniform part of the current flowing on the strip may be represented in the form of the sum of the edge cylindrical waves.
\[\left.\begin{aligned}E_z=-H_\phi&=E_{0z}[f^1(1)e^{ika(\sin\alpha-\sin\phi)}+\\&\qquad+f^1(2)e^{-ika(\sin\alpha-\sin\phi)}]\frac{e^{i(kr+\frac\pi4)}}{\sqrt{2\pi kr}},\\E_\phi=H_z&=H_{0z}[g^1(1)e^{ika(\sin\alpha-\sin\phi)}+\\&\qquad+g^1(2)e^{-ika(\sin\alpha-\sin\phi)}]\frac{e^{i(kr+\frac\pi4)}}{\sqrt{2\pi kr}},\end{aligned}\right\}\tag{6.10}\]
where the functions $f^1$ and $g^1$ are determined in the right half-plane $\left(|\phi|<\frac\pi2\right)$ by the equations
\end{scanpage}
